% =====================================================================
%  CHECKPOINT 01 -- state after block 1 (the skeleton)
%  Copy this over latex-workshop.tex if you need to catch up.
%
%  House style in these files: one sentence per line, \label on the
%  same line as its \section, four spaces per indent.
% =====================================================================

\documentclass[a4paper,11pt]{article}

% ---------------------------------------------------------------------
%  The preamble: everything between \documentclass and \begin{document}.
%  This is where the document is configured. More on it in block 2.
% ---------------------------------------------------------------------
\usepackage[T1]{fontenc}
\usepackage[english]{babel}
\usepackage{newtxtext}

\title{How Fast Does Your Coffee Go Cold?}
\author{Your Name}
\date{\today}

\begin{document}

\maketitle
\tableofcontents

% =====================================================================
\section{Introduction}\label{sec:introduction}

A freshly poured coffee is too hot to drink, and a forgotten one is too cold.
Somewhere in between there is a window of a few minutes in which the drink is actually pleasant,
and this report sets out to find it.
The underlying physics is \emph{Newton's law of cooling},
which says that a body loses heat in proportion to how much hotter it is than its surroundings.
Three \textbf{everyday vessels} are compared,
referred to throughout by the short labels \texttt{mug}, \texttt{glass} and \texttt{flask}.

% TODO 2.2  Add a second paragraph here. Do NOT use \\ to start it --
%           leave an empty line instead and start typing. Say in two
%           sentences why this is worth asking at all (hint: the model
%           has exactly one free parameter per vessel).

% TODO 2.4  After that paragraph, list what the report sets out to do,
%           using an itemize environment. Three bullets is plenty, e.g.
%           stating the cooling model, comparing three vessels, and
%           working out how long each stays drinkable.

% =====================================================================
\section{Theory}\label{sec:theory}

% TODO 2.6  The symbols used in this report need explaining. Build a
%           description list for them:
%             T_0    temperature of the drink at the start
%             T_env  temperature of the room
%             tau    time constant of the vessel
%           Plain text for now -- the maths arrives in block 3.

% =====================================================================
\section{Setup}\label{sec:setup}

% TODO 2.3  The paragraph below was written for you, and it contains
%           three line-break habits that will get you scolded by anyone
%           who reads LaTeX source. Find them and fix them.
%           Reminder: what does \\ actually mean, and where does it
%           legitimately belong?

Each vessel starts with 300 ml of water at 85 degrees Celsius. \\
A thermometer sits in the centre of the liquid, clear of the walls. \\
The temperature is taken every 30 seconds for one hour.

\noindent
The room stays closed and unventilated throughout, at 21 degrees Celsius.

% TODO 2.5  Below, describe the procedure as a numbered list, using an
%           enumerate environment. Four steps, e.g. heat the water,
%           pour it, place the thermometer, record the temperature.

% =====================================================================
\section{Results}\label{sec:results}

% =====================================================================
\section{Conclusion}\label{sec:conclusion}

\end{document}

% =====================================================================
%  BLOCK 2 -- TEXT, LISTS AND THE PREAMBLE
% =====================================================================
%
%  TODO 2.1  First, up in the preamble, add
%
%              \usepackage{parskip}
%
%            Compile before and after and compare the two PDFs. What
%            changed, and why is this nicer than writing \noindent in
%            front of every paragraph?
%
%  TODO 2.2 -- 2.6 are marked in place above, in the order you meet
%            them when reading the document from top to bottom.
%
%  If you finish early
%    * Add \usepackage{microtype} to the preamble. Compile, then look
%      very closely at the right-hand margin.
%    * Nest one list inside another and see what happens to the labels.
% =====================================================================
